\appendix
\section{Reproducibility}
\label{sec:repro}

\paragraph{Frozen campaign.} The 127-case campaign (\texttt{campaign1}) was frozen on 2026-10-20 with the SHA-256 digest of its \texttt{campaign.json} \texttt{\scriptsize 231fea7a9e1722ad46b8c79eeb06789f0437118b8df3fa4976874d7459e4d380}. The file records the hashes of the G0 and G1 Python sources, the native libraries, all benchmark helper files, the registry and the hashed initial atomic densities; every job validates these identities before and after it runs, and each card works in a fork of the campaign whose \texttt{campaign.json} is byte-identical, so the CPU references are shared and one audit applies to all.

\paragraph{Layout.}
The code, benchmark suite and manuscript sources are in the \texttt{improve-single-gpu} branch of \url{https://github.com/qshao/gpu4pyscf} (commit \texttt{ac154f9} at submission). Code and data live in \texttt{benchmarks/single\_gpu/qm\_suite} of the repository (\texttt{cases.json}, \texttt{run.py}, \texttt{schedule.py}, \texttt{adapters.py}, \texttt{audit.py}); the cluster drivers are in \texttt{qm\_suite/lcc}; the stage profiles, storage-tier and capacity runs are in \texttt{benchmarks/single\_gpu/lcc\_profile\_results}; the manuscript is built from \texttt{reports/lcc\_benchmark} (\texttt{collect\_aux.py}, \texttt{make\_figures.py}, \texttt{make\_tables.py}, \texttt{report.tex}). The environment is a conda environment with Python~3.11, PySCF~2.14.0, CuPy~13.4.1, cuTENSOR~2.2.0 and the GCC~12 and CUDA~12.4 toolchains.

\paragraph{Commands.} CPU references: a Slurm job array over the case list calling \texttt{run.py ... --variant CPU --mode reference}. GPU matrix: \texttt{schedule.py --campaign all --cases <list> --repetitions 5} on each card fork. CPU timing: \texttt{run.py ... --variant CPU --mode timing} per case and repetition on whole nodes. Capacity ladder: \texttt{lcc/capacity.py CARD}. Cross-stratum tables: \texttt{lcc/compare\_variants.py}.

\section{Context: earlier single-GPU study on a GB10}
\label{sec:gb10}

Before the cluster campaign, the same branch was compared with the original code on one NVIDIA GB10 (unified memory, weak double-precision throughput, no cuTENSOR bindings so all contractions used the CuPy fallback). Table~\ref{tab:gb10} reproduces that result. The GB10 study motivated two checks in the present report: whether the speed-up survives when cuTENSOR is available (Section~\ref{sec:engine}), and whether the memory reduction matters on cards with far less memory (Section~\ref{sec:memory}).

\begin{table}[!htbp]
\centering
\small
\resizebox{\textwidth}{!}{%
\begin{tabular}{lrrrrrrr}
\toprule
system & atoms & AOs & original (s) & new (s) & ratio & orig. GB & new GB \\
\midrule
Vitamin C & 20 & 208 & 5.49 & 5.28 & 1.04 & 1.4 & 0.5 \\
Tamoxifen & 57 & 537 & 70.5 & 48.1 & 1.47 & 21.1 & 4.3 \\
Taxol & 113 & 1123 & 1103 & 566.0 & 1.95 & 65.6 & 30.1 \\
\bottomrule
\end{tabular}}
\caption{Earlier study on one NVIDIA GB10 (unified memory, CuPy fallback contractions, B3LYP/def2-SVP, three timing repetitions); shown for context. The original code fails the CPU gradient gate on Taxol, so that ratio is not an equal-accuracy speed-up.}
\label{tab:gb10}
\end{table}


\begin{landscape}
\section{Per-case results}
\scriptsize\setlength{\tabcolsep}{3pt}
\begin{longtable}{lrrrrrrrrrrr}
\caption{Per-case results. Columns A100 to P100-12: time ratio original/new (``nc'': a solver or reference did not converge; ``mem'': resource limit; ``acc'': the original code fails the accuracy gate; ``--'': not run or not qualified). G1: new-code time on the V100; mem: V100 GPU memory original/new in GB; CPU: best CPU time.}\label{tab:percase}\\
\toprule
case & AOs & A100 & V100 & A100c & V100c & P100-16 & P100-12 & G1 V100 (s) & mem (GB) & CPU (s) \\
\midrule
\endfirsthead
\toprule case & AOs & A100 & V100 & A100c & V100c & P100-16 & P100-12 & G1 V100 (s) & mem (GB) & CPU (s) \\ \midrule \endhead
\bottomrule \endfoot
\texttt{c-acetate-b3lyp-svp} & 110 & 0.99 & 1.01 & 1.01 & 1.00 & 1.00 & 1.03 & 1.28 & 0.1/0.1 & 2.18 \\
\texttt{c-acetate-pbe-svp} & 110 & 1.00 & 0.97 & 1.00 & 0.98 & 0.99 & 1.01 & 1.37 & 0.1/0.1 & 2.45 \\
\texttt{c-benzene-b3lyp-svp} & 114 & 0.99 & 0.98 & 1.01 & 1.01 & 1.00 & 0.99 & 1.31 & 0.2/0.2 & 3.67 \\
\texttt{c-benzene-b3lyp-tzvp} & 222 & 1.01 & 1.02 & 1.02 & 1.01 & 1.03 & 1.07 & 1.59 & 0.9/0.3 & 6.26 \\
\texttt{c-benzene-hf-svp} & 114 & 1.02 & 1.03 & 1.02 & 1.04 & 1.01 & 1.06 & 0.77 & 0.2/0.1 & 0.86 \\
\texttt{c-benzene-lda-svp} & 114 & 0.99 & 0.98 & 1.02 & 1.01 & 1.00 & 1.01 & 1.20 & 0.1/0.1 & 2.79 \\
\texttt{c-benzene-pbe-svp} & 114 & 1.02 & 0.99 & 1.04 & 1.06 & 1.00 & 1.00 & 1.29 & 0.1/0.1 & 3.80 \\
\texttt{c-benzene-pbe-tzvp} & 222 & 1.00 & 0.99 & 1.00 & 1.01 & 0.99 & 1.02 & 1.55 & 0.2/0.2 & 5.96 \\
\texttt{c-benzene-tpss-svp} & 114 & 1.00 & 1.01 & 1.15 & 1.14 & 1.18 & 1.21 & 1.61 & 0.1/0.1 & 4.74 \\
\texttt{c-benzene-wb97x-svp} & 114 & 1.00 & 0.96 & 1.06 & 0.99 & 1.00 & 1.02 & 1.74 & 0.2/0.2 & 4.00 \\
\texttt{c-benzene\_dimer-b3lyp-svp} & 228 & 0.98 & 0.98 & 0.99 & 1.01 & 1.01 & 1.01 & 2.08 & 1.5/0.7 & 19.9 \\
\texttt{c-benzene\_dimer-pbe-svp} & 228 & 1.02 & 0.97 & 1.01 & 0.99 & 1.02 & 1.00 & 1.96 & 0.5/0.4 & 18.1 \\
\texttt{c-distorted\_h2o-b3lyp-svp} & 24 & 0.98 & 1.00 & 1.02 & 1.00 & 1.01 & 1.00 & 1.03 & 0.0/0.0 & 0.51 \\
\texttt{c-h2o-b3lyp-pvdz} & 24 & 1.01 & 0.95 & 1.00 & 0.98 & 0.99 & 0.97 & 1.06 & 0.0/0.0 & 0.49 \\
\texttt{c-h2o-b3lyp-svp} & 24 & 1.01 & 1.04 & 1.01 & 1.00 & 0.98 & 1.04 & 0.99 & 0.0/0.0 & 0.51 \\
\texttt{c-h2o-hf-svp} & 24 & 1.04 & 1.04 & 1.03 & 1.01 & 0.99 & 1.01 & 0.71 & 0.0/0.0 & 0.38 \\
\texttt{c-h2o-lda-svp} & 24 & 1.01 & 1.03 & 0.98 & 0.95 & 1.00 & 1.05 & 1.00 & 0.0/0.0 & 0.56 \\
\texttt{c-h2o-pbe-pvdz} & 24 & 1.01 & 1.01 & 1.01 & 1.03 & 0.97 & 1.00 & 0.98 & 0.0/0.0 & 0.60 \\
\texttt{c-h2o-pbe-svp} & 24 & 1.00 & 1.03 & 0.99 & 1.01 & 1.00 & 1.08 & 0.97 & 0.0/0.0 & 0.59 \\
\texttt{c-h2o-tpss-svp} & 24 & 1.02 & 1.01 & 1.08 & 1.05 & 1.07 & 1.16 & 1.02 & 0.0/0.0 & 0.67 \\
\texttt{c-h2o-wb97x-svp} & 24 & 0.99 & 0.98 & 0.99 & 1.01 & 1.01 & 1.02 & 1.16 & 0.0/0.0 & 0.52 \\
\texttt{c-inosine-b3lyp-svp} & 326 & 1.01 & 0.99 & 1.02 & 1.05 & 1.05 & 1.09 & 3.98 & 4.6/1.3 & 72.4 \\
\texttt{c-inosine-hf-svp} & 326 & 0.94 & 0.95 & 0.95 & 1.03 & 1.06 & 1.13 & 2.00 & 4.6/1.3 & 17.3 \\
\texttt{c-inosine-lda-svp} & 326 & 1.01 & 0.99 & 1.01 & 1.00 & 1.00 & 1.00 & 3.11 & 0.8/0.6 & 48.5 \\
\texttt{c-inosine-pbe-svp} & 326 & 1.03 & 1.02 & 1.03 & 1.02 & 1.02 & 1.03 & 3.81 & 0.8/0.6 & 68.2 \\
\texttt{c-inosine-tpss-svp} & 326 & 1.00 & 1.00 & 1.24 & 1.24 & 1.37 & 1.53 & 5.22 & 0.8/0.6 & 106.8 \\
\texttt{c-inosine-wb97x-svp} & 326 & 0.99 & 0.98 & 0.99 & 1.02 & 1.04 & 1.08 & 4.97 & 4.6/1.5 & 77.3 \\
\texttt{c-methylammonium-b3lyp-svp} & 58 & 1.02 & 0.98 & 1.01 & 0.98 & 1.00 & 1.04 & 1.19 & 0.1/0.1 & 1.37 \\
\texttt{c-methylammonium-pbe-svp} & 58 & 0.99 & 0.98 & 0.99 & 1.01 & 1.01 & 1.01 & 1.22 & 0.1/0.1 & 1.57 \\
\texttt{c-no-b3lyp-svp} & 28 & 1.01 & 0.98 & acc & acc & 1.02 & 0.98 & 1.11 & 0.0/0.0 & 0.63 \\
\texttt{c-no-pbe-svp} & 28 & nc & nc & nc & 0.98 & -- & nc & 1.07 & 0.0/0.0 & -- \\
\texttt{c-o2-b3lyp-svp} & 28 & 0.98 & 1.01 & 1.00 & 1.00 & 1.01 & 0.97 & 1.01 & 0.0/0.0 & 0.56 \\
\texttt{c-o2-pbe-svp} & 28 & 1.01 & 1.02 & 1.02 & 0.99 & 1.00 & 0.97 & 1.01 & 0.0/0.0 & 0.63 \\
\texttt{c-oh-b3lyp-svp} & 19 & -- & -- & -- & nc & -- & -- & -- & -- & -- \\
\texttt{c-oh-pbe-svp} & 19 & -- & -- & -- & -- & nc & -- & -- & -- & -- \\
\texttt{c-stretched\_n2-b3lyp-svp} & 28 & 1.00 & 0.98 & 0.97 & 1.00 & 1.03 & 1.00 & 0.96 & 0.0/0.0 & 0.47 \\
\texttt{c-tamoxifen-b3lyp-svp} & 537 & 0.96 & 0.96 & 1.31 & 1.34 & 1.39 & 1.77 & 7.80 & 15.9/4.3 & 310.6 \\
\texttt{c-tamoxifen-b3lyp-tzvp} & 1042 & 0.92 & 0.96 & 2.39 & 1.76 & -- & -- & 18.2 & 24.3/11.6 & 593.6 \\
\texttt{c-tamoxifen-hf-svp} & 537 & 0.88 & 0.91 & 2.34 & 1.75 & 1.63 & 2.26 & 4.61 & 15.8/4.3 & 82.2 \\
\texttt{c-tamoxifen-lda-svp} & 537 & 1.00 & 0.99 & 1.01 & 1.00 & 1.00 & 1.00 & 4.81 & 2.4/0.6 & 201.4 \\
\texttt{c-tamoxifen-pbe-svp} & 537 & 1.01 & 1.01 & 1.03 & 1.02 & 1.02 & 1.04 & 5.70 & 2.5/0.7 & 251.4 \\
\texttt{c-tamoxifen-pbe-tzvp} & 1042 & 1.01 & 1.00 & 1.00 & 1.00 & 1.01 & 1.01 & 10.8 & 2.6/1.0 & 392.7 \\
\texttt{c-tamoxifen-tpss-svp} & 537 & 1.01 & 1.00 & 1.24 & 1.24 & 1.36 & 1.43 & 8.72 & 2.5/0.8 & 381.9 \\
\texttt{c-tamoxifen-wb97x-svp} & 537 & 0.94 & 0.94 & 1.48 & 1.27 & 1.32 & 1.64 & 10.8 & 15.9/5.2 & 371.2 \\
\texttt{c-taxol-b3lyp-svp} & 1123 & acc & acc & acc & -- & mem & mem & 58.8 & 29.3/20.0 & 2472 \\
\texttt{c-taxol-hf-svp} & 1123 & acc & acc & acc & -- & mem & mem & 47.5 & 29.5/20.0 & 998.5 \\
\texttt{c-taxol-lda-svp} & 1123 & 1.01 & 1.00 & 1.00 & 1.00 & -- & 1.01 & 18.3 & 4.9/1.0 & 1377 \\
\texttt{c-taxol-pbe-svp} & 1123 & 1.03 & 1.09 & 1.03 & 1.09 & -- & 1.09 & 24.2 & 5.1/1.4 & 1604 \\
\texttt{c-taxol-tpss-svp} & 1123 & 1.00 & 1.00 & 1.20 & 1.05 & -- & 1.21 & 38.1 & 5.1/1.4 & 2063 \\
\texttt{c-taxol-wb97x-svp} & 1123 & acc & acc & acc & -- & mem & mem & 85.1 & 29.6/25.9 & 3006 \\
\texttt{c-vitamin\_c-b3lyp-svp} & 208 & 1.04 & 1.01 & 1.03 & 1.02 & 1.02 & 1.03 & 2.18 & 1.2/0.5 & 17.3 \\
\texttt{c-vitamin\_c-b3lyp-tzvp} & 420 & 1.00 & 0.99 & 1.09 & 1.14 & 1.12 & 1.19 & 3.24 & 5.7/1.5 & 34.0 \\
\texttt{c-vitamin\_c-hf-svp} & 208 & 0.89 & 0.99 & 1.02 & 1.00 & 1.03 & 1.02 & 1.08 & 1.2/0.4 & 4.62 \\
\texttt{c-vitamin\_c-lda-svp} & 208 & 0.99 & 0.99 & 1.00 & 0.99 & 0.99 & 1.00 & 2.03 & 0.3/0.3 & 12.3 \\
\texttt{c-vitamin\_c-pbe-svp} & 208 & 1.03 & 1.01 & 1.01 & 1.00 & 1.03 & 0.98 & 2.33 & 0.4/0.3 & 16.8 \\
\texttt{c-vitamin\_c-pbe-tzvp} & 420 & 1.00 & 1.01 & 1.00 & 1.01 & 1.00 & 0.99 & 3.19 & 0.4/0.4 & 32.7 \\
\texttt{c-vitamin\_c-tpss-svp} & 208 & 1.00 & 0.99 & 1.22 & 1.22 & 1.33 & 1.54 & 3.14 & 0.4/0.3 & 20.9 \\
\texttt{c-vitamin\_c-wb97x-svp} & 208 & 1.01 & 1.01 & 1.02 & 1.00 & 1.02 & 1.02 & 2.66 & 1.2/0.6 & 19.4 \\
\texttt{c-water\_dimer-b3lyp-svp} & 48 & 0.99 & 1.01 & 1.00 & 1.01 & 0.99 & 0.99 & 1.12 & 0.1/0.1 & 1.05 \\
\texttt{c-water\_dimer-pbe-svp} & 48 & 1.00 & 1.02 & 1.00 & 0.99 & 1.02 & 1.05 & 1.16 & 0.1/0.1 & 1.21 \\
\texttt{direct-benzene-b3lyp} & 114 & 1.00 & 1.00 & 0.99 & 0.99 & 1.00 & 1.00 & 2.13 & 0.2/0.2 & 6.45 \\
\texttt{direct-benzene-hf} & 114 & 1.00 & 1.01 & 0.99 & 0.99 & 0.99 & 0.98 & 1.61 & 0.2/0.2 & 3.66 \\
\texttt{direct-benzene-pbe} & 114 & 1.00 & 1.03 & 1.00 & 0.98 & 0.99 & 1.02 & 1.62 & 0.2/0.2 & 6.30 \\
\texttt{direct-h2o-b3lyp} & 24 & 1.00 & 1.02 & 1.00 & 0.98 & 1.01 & 0.98 & 1.41 & 0.2/0.2 & 0.55 \\
\texttt{direct-h2o-hf} & 24 & 1.00 & 0.99 & 1.00 & 1.03 & 0.97 & 1.02 & 1.10 & 0.2/0.2 & 0.38 \\
\texttt{direct-h2o-pbe} & 24 & 0.99 & 1.05 & 0.98 & 1.01 & 1.01 & 1.01 & 0.97 & 0.2/0.2 & 0.50 \\
\texttt{e-benzene-b3lyp-hessian} & 114 & 1.00 & 1.00 & 1.04 & 1.03 & 1.04 & 1.02 & 8.07 & 0.4/0.4 & 699.3 \\
\texttt{e-benzene-b3lyp-tda} & 114 & nc & nc & nc & nc & nc & nc & -- & 0.4/0.4 & 66.7 \\
\texttt{e-benzene-b3lyp-tddft} & 114 & nc & nc & nc & nc & nc & nc & -- & 0.9/0.9 & 32.9 \\
\texttt{e-benzene-dfmp2} & 114 & 0.96 & 1.02 & 1.02 & 1.03 & 1.08 & 1.02 & 0.75 & 0.2/0.2 & 0.58 \\
\texttt{e-benzene-hf-hessian} & 114 & 1.01 & 1.02 & 1.04 & 1.01 & 1.04 & 1.01 & 1.95 & 0.3/0.3 & 140.7 \\
\texttt{e-benzene-optimization} & 114 & 1.00 & 1.01 & 1.02 & 1.05 & 1.02 & 1.01 & 5.15 & 0.2/0.2 & 19.2 \\
\texttt{e-benzene-pbe-hessian} & 114 & 1.00 & 1.01 & 1.02 & 1.03 & 1.04 & 1.03 & 7.71 & 0.5/0.5 & 816.1 \\
\texttt{e-benzene-polarizability} & 114 & 1.00 & 1.01 & 0.99 & 0.98 & 1.01 & 0.99 & 2.90 & 0.1/0.1 & 3.94 \\
\texttt{e-benzene-shielding} & 114 & 1.00 & 0.99 & 1.00 & 1.00 & 1.01 & 1.01 & 7.51 & 0.2/0.2 & 5.52 \\
\texttt{e-benzene-wb97mv} & 114 & acc & acc & acc & acc & acc & acc & -- & 0.2/0.2 & 55.6 \\
\texttt{e-benzene-wb97x-tda} & 114 & nc & nc & nc & nc & nc & nc & -- & 0.4/0.4 & 159.9 \\
\texttt{e-benzene-wb97x-tddft} & 114 & nc & nc & nc & nc & nc & nc & -- & 0.9/0.9 & 62.7 \\
\texttt{e-ch3i-pbe-ecp} & 55 & 1.02 & 0.99 & 1.01 & 1.02 & 1.01 & 1.01 & 1.80 & 0.1/0.1 & 2.38 \\
\texttt{e-formaldehyde-b3lyp-tda} & 38 & 1.01 & 1.01 & 1.00 & 1.00 & 1.01 & 0.99 & 2.91 & 0.1/0.1 & 1.83 \\
\texttt{e-formaldehyde-b3lyp-tddft} & 38 & 0.99 & 1.00 & 1.01 & 0.99 & 0.99 & 1.01 & 1.52 & 0.1/0.1 & 2.46 \\
\texttt{e-formaldehyde-ccsd} & 38 & 0.91 & 1.00 & 1.00 & 0.91 & mem & mem & 8.38 & 0.1/0.1 & 1.22 \\
\texttt{e-formaldehyde-optimization} & 38 & 1.00 & 1.01 & 1.01 & 1.02 & 1.00 & 1.02 & 4.23 & 0.1/0.1 & 3.18 \\
\texttt{e-formaldehyde-wb97x-tda} & 38 & 0.99 & 0.99 & 0.99 & 1.00 & 1.01 & 1.01 & 10.9 & 0.1/0.1 & 1.96 \\
\texttt{e-formaldehyde-wb97x-tddft} & 38 & nc & nc & nc & nc & nc & nc & -- & 0.1/0.1 & 2.74 \\
\texttt{e-h2o-b3lyp-d3bj} & 24 & 1.02 & 0.97 & 1.01 & 1.04 & 1.02 & 0.92 & 0.98 & 0.0/0.0 & 0.49 \\
\texttt{e-h2o-b3lyp-d4} & 24 & 0.98 & 1.00 & 0.99 & 0.98 & 1.01 & 0.96 & 1.00 & 0.0/0.0 & 0.54 \\
\texttt{e-h2o-b3lyp-hessian} & 24 & 1.00 & 0.98 & 1.04 & 1.02 & 0.99 & 1.00 & 2.24 & 0.1/0.1 & 9.57 \\
\texttt{e-h2o-b3lyp-iefpcm} & 24 & 1.00 & 1.00 & 1.00 & 0.97 & 1.00 & 1.03 & 1.25 & 0.0/0.0 & 0.79 \\
\texttt{e-h2o-b3lyp-smd} & 24 & 0.99 & 1.03 & 1.00 & 0.95 & 0.97 & 1.01 & 1.23 & 0.0/0.0 & 0.81 \\
\texttt{e-h2o-ccsd} & 24 & 1.00 & 1.00 & 1.01 & 1.01 & mem & mem & 6.95 & 0.1/0.1 & 0.73 \\
\texttt{e-h2o-dfmp2} & 24 & 0.98 & 1.01 & 1.01 & 1.04 & 1.01 & 0.90 & 0.64 & 0.0/0.0 & 0.34 \\
\texttt{e-h2o-hf-hessian} & 24 & 1.00 & 1.01 & 1.01 & 1.01 & 1.02 & 1.04 & 1.29 & 0.0/0.0 & 4.14 \\
\texttt{e-h2o-optimization} & 24 & 1.00 & 1.02 & 1.00 & 1.00 & 1.01 & 1.02 & 2.74 & 0.0/0.0 & 1.40 \\
\texttt{e-h2o-pbe-hessian} & 24 & 1.01 & 1.03 & 1.01 & 1.02 & 1.01 & 1.00 & 2.14 & 0.1/0.1 & 15.3 \\
\texttt{e-h2o-polarizability} & 24 & 1.00 & 1.01 & 1.01 & 0.99 & 0.99 & 0.95 & 1.65 & 0.1/0.1 & 0.39 \\
\texttt{e-h2o-shielding} & 24 & 1.00 & 1.01 & 1.00 & 1.02 & 1.01 & 0.99 & 1.74 & 0.1/0.1 & 0.41 \\
\texttt{e-h2o-wb97mv} & 24 & 1.00 & 1.00 & 1.03 & 1.05 & 1.06 & 0.94 & 1.56 & 0.0/0.0 & 3.69 \\
\texttt{e-i2-pbe-ecp} & 52 & 1.00 & 1.00 & 1.02 & 0.99 & 1.00 & 1.00 & 1.38 & 0.1/0.1 & 2.42 \\
\texttt{e-inosine-b3lyp-iefpcm} & 326 & 0.99 & 0.99 & 1.01 & 1.03 & 1.04 & 1.06 & 5.30 & 4.9/1.9 & 186.7 \\
\texttt{e-inosine-b3lyp-smd} & 326 & 1.00 & 0.99 & 1.01 & 1.03 & 1.04 & 1.07 & 5.26 & 5.0/1.6 & 199.2 \\
\texttt{e-inosine-b3lyp-tda} & 326 & 1.00 & 1.01 & 1.01 & 1.01 & 1.01 & 1.02 & 131.5 & 1.7/1.7 & 688.8 \\
\texttt{e-inosine-b3lyp-tddft} & 326 & 1.00 & 1.01 & 1.01 & 1.03 & 1.07 & 1.05 & 9.87 & 5.6/5.6 & 790.2 \\
\texttt{e-inosine-wb97x-tda} & 326 & nc & nc & nc & nc & nc & nc & -- & 2.1/2.1 & 889.1 \\
\texttt{e-inosine-wb97x-tddft} & 326 & 1.00 & 1.00 & 1.01 & 1.03 & 1.06 & 1.05 & 15.6 & 5.8/5.8 & 1421 \\
\texttt{e-no-dfump2} & 28 & 1.00 & 1.01 & 1.03 & 1.03 & 1.00 & 0.94 & 0.71 & 0.0/0.0 & 0.43 \\
\texttt{e-oh-dfump2} & 19 & 0.97 & 0.99 & 0.99 & 1.03 & 1.01 & 1.07 & 0.69 & 0.0/0.0 & 0.38 \\
\texttt{e-tamoxifen-b3lyp-d3bj} & 537 & 0.96 & 0.96 & 1.31 & 1.34 & 1.38 & 1.81 & 7.80 & 15.9/4.3 & 308.9 \\
\texttt{e-tamoxifen-b3lyp-d4} & 537 & 0.95 & 0.96 & 1.31 & 1.34 & 1.37 & 1.77 & 7.73 & 15.9/4.3 & 310.3 \\
\texttt{e-tamoxifen-b3lyp-iefpcm} & 537 & 0.97 & 0.97 & 1.20 & 1.23 & 1.26 & 1.52 & 10.4 & 16.4/5.7 & 927.4 \\
\texttt{e-tamoxifen-b3lyp-smd} & 537 & 0.97 & 0.98 & 1.20 & 1.23 & 1.25 & 1.53 & 10.4 & 16.5/4.5 & 954.7 \\
\texttt{e-tamoxifen-wb97mv} & 537 & acc & acc & acc & acc & acc & acc & -- & 15.9/5.6 & 2768 \\
\texttt{e-vitamin\_c-b3lyp-hessian} & 208 & 1.00 & 0.99 & 1.04 & 1.04 & 1.04 & 1.04 & 20.9 & 2.6/2.6 & 5268 \\
\texttt{e-vitamin\_c-b3lyp-iefpcm} & 208 & 1.02 & 1.00 & 1.02 & 1.01 & 1.02 & 1.02 & 2.91 & 1.4/0.7 & 46.2 \\
\texttt{e-vitamin\_c-b3lyp-smd} & 208 & 1.01 & 1.00 & 1.02 & 1.01 & 1.01 & 1.04 & 2.93 & 1.4/0.7 & 50.3 \\
\texttt{e-vitamin\_c-dfmp2} & 208 & 1.01 & 0.97 & 1.06 & 1.03 & 1.03 & 0.91 & 1.17 & 0.8/0.8 & 2.13 \\
\texttt{e-vitamin\_c-optimization} & 208 & 1.00 & 1.00 & 1.03 & 1.02 & 1.03 & 1.04 & 48.6 & 1.2/0.5 & 479.3 \\
\texttt{e-vitamin\_c-pbe-hessian} & 208 & 1.01 & 1.00 & 1.03 & 1.04 & 1.02 & 1.04 & 18.7 & 2.8/2.8 & 6185 \\
\texttt{e-vitamin\_c-polarizability} & 208 & 1.00 & 1.01 & 1.01 & 1.00 & 1.00 & 1.00 & 7.21 & 0.2/0.2 & 20.6 \\
\texttt{e-vitamin\_c-shielding} & 208 & 1.00 & 1.00 & 1.00 & 1.01 & 1.01 & 1.01 & 28.0 & 0.3/0.3 & 19.1 \\
\texttt{e-vitamin\_c-wb97mv} & 208 & acc & acc & acc & acc & acc & acc & -- & 1.3/0.7 & 239.3 \\
\texttt{e-water\_dimer-dfmp2} & 48 & 0.98 & 1.01 & 1.00 & 1.04 & 1.00 & 0.95 & 0.67 & 0.0/0.0 & 0.36 \\
\texttt{pbc-mgo-gamma} & 27 & acc & acc & acc & acc & acc & acc & -- & 0.0/0.0 & 3.57 \\
\texttt{pbc-si-gamma} & 26 & 0.97 & 1.00 & 1.00 & 1.02 & 1.01 & 1.00 & 0.94 & 0.0/0.0 & 2.92 \\
\texttt{pbc-si-k222} & 26 & 0.98 & 1.05 & 1.01 & 1.03 & 1.04 & 1.04 & 1.40 & 0.2/0.2 & 7.84 \\
\texttt{pbc-si-supercell-gamma} & 208 & 1.00 & 1.02 & 1.04 & 1.03 & 1.00 & 0.99 & 2.12 & 0.1/0.1 & 58.6 \\
\texttt{probe-oh-wb97x-svp} & 19 & -- & -- & -- & -- & -- & -- & -- & -- & -- \\
\end{longtable}

\end{landscape}
